% ==== note-gen Chinese preamble (中英夹杂 notes). Compile with XeLaTeX. ====
\documentclass[11pt,a4paper]{article}
\usepackage[a4paper,top=2.5cm,bottom=2.5cm,left=2.5cm,right=2.5cm]{geometry}
\usepackage{ctex}                 % Chinese typesetting (requires xelatex)
\usepackage{amsmath,amssymb,amsthm}
\usepackage{graphicx}
\usepackage{float}
\usepackage{booktabs}
\usepackage{array}
\usepackage{xcolor}
\usepackage{tcolorbox}
\usepackage{listings}
\usepackage[colorlinks=true,allcolors=blue]{hyperref}
\usepackage{enumitem}
\usepackage{setspace}
\usepackage{cite}

\onehalfspacing
\newtheorem{definition}{定义}
\newtheorem{intuition}{直觉}

% Highlighted callout boxes (English titles, 中文 content)
\newtcolorbox{keybox}[1][]{colback=blue!5!white,colframe=blue!60!black,
  fonttitle=\bfseries,title=Key Idea,#1}
\newtcolorbox{pitfall}[1][]{colback=red!5!white,colframe=red!60!black,
  fonttitle=\bfseries,title=Common Pitfall,#1}

\lstset{
  basicstyle=\ttfamily\small,
  backgroundcolor=\color{gray!8},
  frame=single,
  rulecolor=\color{gray!35},
  breaklines=true,
  columns=fullflexible,
  keepspaces=true,
  showstringspaces=false
}

\newcommand{\eg}{\textit{e.g.}}
\newcommand{\ie}{\textit{i.e.}}
